\scanpage{18}
% archive sheet number: 17
% L1--3.
\diagram{\begin{tikzcd}[ampersand replacement=\&,column sep=large]
\prod J_x^{(0)}\arrow[r]\arrow[d,equal] \& J_\xi^*\\
J_1 \& J_0=J_{K/k}
\end{tikzcd}}
\unsure{ou}
\[
\deleted{\unsure{\mathrm{gr}}}\quad J_{C/k}^*\quad\text{jacobienne de }C
\]
% L4--12: crossed block. All its readable content is retained as deleted.
\deleted{La chose \unsure{spéciale|précise} et surprenante, c'est que l'on a aussi $J_{C/k}^*\simeq\underline{H}^1(C,G_m)$, \unsure{et} [si on \unsure{prend} le complexe $\operatorname{Hom}_{k\text{-}\unsure{\mathrm{gr}}}(\unsure{J_*|J^*},G_m)$, on trouve}
\[
\deleted{K^*\longrightarrow\substack{\text{diviseurs}\\\text{\unsure{ordinaires}}}}
\]
\deleted{] et que $J_x^{(0)}(k)=G_m(\unsure{\hat{\mathcal O}_x|\hat O_x})$.}

\deleted{Il faut aussi expliquer le rôle \unsure{du|de} $J_x^*$.}
% L13--19.
La chose \unsure{non|un} \unsure{primordiale|principale} : \unsure{primo}, c'est que le complexe
\[
J_1\unsure{\longrightarrow|\text{---}}J_0
\]
(ou plutôt $J_1(k)\unsure{\longrightarrow|\text{---}}J_0(k)$) \unsure{n'est} autre que $H_{\mathrm{II}}DD(G_m)$, on peut donc, \unsure{moyennant} un abus de \unsure{notations},, écrire
\[
H_{\mathrm I}DD(G)=\operatorname{Hom}\bigl(\underline{H_{\mathrm{II}}DD}(G_m),G\bigr)
\]
\deleted{ce qui suggère que l'on ait}
